% 8.模型改进推广.tex —— 八、模型的改进与推广
\section{模型的改进与推广}

\subsection{模型的改进}
后续可计入蒸发潜热汇项 $-\rho L_v\,\partial C/\partial t$，使热质真正双向耦合，并以联立迭代求解；水分边界条件若题意采用 Dirichlet 型（表面浓度恒等于平衡含水率），应替换式 \eqref{eq:robinC} 重算。此外可加密网格至 $N=160$ 使温度误差降至 $2.6\times10^{-5}$ $^\circ$C。

\subsection{模型的推广}
本模型的“径向一维 + 节点分类显式/隐式差分 + Thomas 追赶”框架可推广到其他具有柱/球对称性的瞬态扩散问题（如传热、传质、充电/放电扩散过程）。推广时只需替换扩散系数与边界条件，中心点、内部点、边界点的三类离散结构与求解流程保持不变。
